\chapter{Methodology and Implementation}

\section{Methodology}

\subsection{Research Methodology}
Explain:

prototype-oriented approach,
modular architecture,
experimental evaluation.

\subsection{System Requirements Analysis}
Functional Requirements
Non-functional Requirements

\subsection{Proposed System Architecture}

Introduce:

layered architecture,
overall system design.

Insert:

overall architecture diagram.

\subsection{Proposed Workflow}

Explain:

authentication,
report creation,
moderation,
blockchain anchoring.

Insert:

workflow diagram.

\subsection{Project Plan and Timeline}

Include:

Gantt chart,
implementation phases.

\subsection{Expected Outcomes}

Examples:

anonymous reporting,
immutable records,
decentralized moderation,
transparency.

\section{Implementation}


\subsection{Permissioned Blockchain Implementation}
Geth PoA Network Setup
Validator Node Configuration
RPC Communication

\subsection{Smart Contract Implementation}
Report Contract
Ticket Validation
State Transition Logic

\subsection{Identity and Access Control Implementation}
GovID Simulator
Ticket Batch Generation
Citizen Keypair Generation
\subsection{IPFS Off-chain Storage Integration}

The system employs IPFS as the off-chain storage layer dedicated to managing citizen-submitted report images. Rather than recording image files directly on the blockchain, only the Content Identifier (CID) produced by IPFS is passed to the smart contract for permanent on-chain recording. The CID is a cryptographic hash that uniquely represents the exact content of a file, such that any modification to the image produces an entirely different CID. This content-addressing mechanism ensures tamper-evident image storage while significantly reducing on-chain data overhead and improving overall system scalability.

The IPFS integration is implemented as a dedicated microservice structured using a layered architecture comprising route definitions, controller components, utility modules, and middleware layers. The middleware layer enforces file validation by applying image type and size restrictions before any content reaches the storage layer. The controller layer manages the core operations of image upload, CID-based image retrieval, CID verification, and pin management. The utility layer handles direct communication with the IPFS node through the Kubo RPC client.

The image upload workflow is managed entirely by the backend server. When a citizen submits a report through the DApp, the report data and attached image are sent to the backend server. The backend server forwards the image to the AI moderation service for content analysis. If the image passes moderation, the backend server sends only the approved image to the IPFS microservice. The IPFS microservice uploads the image to the IPFS node, which pins the content permanently to prevent garbage collection, and returns the resulting CID to the microservice. The microservice passes this CID back to the backend server, which then combines the image CID with the report data and submits both to the smart contract for permanent on-chain recording.

During retrieval, the image CID is read from the blockchain by the backend server and forwarded to the IPFS microservice, which fetches the corresponding image from the node and returns it for display in the DApp. Since the CID is deterministically derived from the image content, the retrieved image can be independently verified against the on-chain record, confirming that the evidence has not been altered since its original submission.

The service additionally incorporates a CID verification mechanism that confirms a CID is resolvable on the IPFS network before the backend server submits it to the smart contract, preventing broken references from being recorded on-chain. A pin management mechanism is also included to handle image removal where necessary, such as when content is later identified as violating system guidelines. Removing the pin from the local node allows garbage collection to reclaim storage while the on-chain CID record is preserved as an immutable audit trail, directly addressing the toxic immutability challenge identified in the literature.

The complete service is containerised using Docker and deployed alongside the IPFS node as a two-container stack, with the IPFS microservice exposed externally on port 4000 and the IPFS node operating internally within the Docker network.

At the current stage of development, the system operates with a single IPFS node within the deployment environment. As a planned future implementation, three IPFS nodes will be deployed across three separate machines and connected as peers within a private controlled network. In this configuration, when an image is uploaded and pinned on the primary node, the two peer nodes will retrieve and independently pin the same content, ensuring redundant availability across all three machines. This multi-node setup will demonstrate the fault tolerance and decentralised storage properties of IPFS at prototype scale, and will operate entirely within the controlled laboratory environment without connecting to the public IPFS network, consistent with the defined project scope.

\subsection{AI-based Moderation System Implementation}


Balanced coverage only.

Subsections:

Oracle Architecture
Oracle Workers
Aggregation Logic
Moderation Workflow
Security Mechanisms

Do NOT over-expand this section.

\subsection{Web-based DApp Development}
User Authentication Interface
Report Submission Interface
Blockchain Query Interface
\subsection{Deployment Environment}

Very important.

Discuss:

VPS infrastructure,
Ubuntu server,
Docker containers,
service distribution.

Include deployment diagram.
